% !PRINTED 442
\weberpage{442}
%<<CONT
\[
4.\quad \left(\frac{-1}{n}\right) = (-1)^{\frac{n-1}{2}}.
\]

Setzen wir $2m$ für $m$ und benutzen die Formel

\[
\sin \frac{4\nu m \pi}{n} = 2 \sin \frac{2\nu m \pi}{n} \cos \frac{2\nu m \pi}{n},
\]

so folgt aus (3)

\[
5.\quad \left(\frac{2m}{n}\right) = (-1)^{\frac{n^2-1}{8}} \left(\frac{m}{n}\right)
\]

\gesperrt{und für} $m = 1$

\[
6.\quad \left(\frac{2}{n}\right) = (-1)^{\frac{n^2-1}{8}}.
\]

Aendert man $m$ um ein Vielfaches von $n$, so bleiben alle
Factoren des Productes (4) ungeändert, und daraus folgt:

7. \gesperrt{Es ist}

\[
\left(\frac{m}{n}\right) = \left(\frac{m'}{n}\right),
\]

\gesperrt{wenn} $m \equiv m'$ \gesperrt{(mod} $n$\gesperrt{).}

Wir setzen jetzt nicht nur $n$, sondern auch $m$ als ungerade
und positiv voraus und wenden die Formel (29) des \S.\,137 an,
in der wir $n$ durch $m$ ersetzen:

\weq{5}{\frac{\sin m \varphi}{\sin \varphi} = (-1)^{\frac{m-1}{2}} 2^{m-1} \overset{\beta}{\Pi} (\sin^2 \varphi - \beta),}

worin $\beta$ die sämmtlichen Wurzeln der Gleichung $\Phi_m(x) = 0$
durchläuft, die den Ausdruck

\weq{6}{\beta = \sin^2 \frac{2\mu\pi}{m}}

% [corrected 2026-08-07, 10x zoom strip 7/10] The page prints "haben." -- the fragment had
% "haben wir.", a fabricated word. The relative clause is "die den Ausdruck (6) haben", which
% "wir" makes ungrammatical, so the fragment's own German was evidence before the crop was.
% First plain-prose ADDITION found in this volume: the audit lanes had been finding omissions,
% and an inserted word is the harder direction to notice, because nothing looks missing.
haben. Ist nun $m'$ eine zu $n$ theilerfremde Zahl, so durchläuft

\weq{7}{\alpha = \sin^2 \frac{2\nu m' \pi}{n}}

die sämmtlichen Wurzeln der Gleichung $\Phi_n(x) = 0$, und wenn
wir in (5)

\weq{8}{\varphi = \frac{2\nu m' \pi}{n}}

setzen und das Product in Bezug auf alle $\nu$ nehmen, so folgt:

\weq{9}{\frac{\overset{\nu}{\Pi} \sin \dfrac{2\nu m m' \pi}{n}}{\overset{\nu}{\Pi} \sin \dfrac{2\nu m' \pi}{n}} = (-1)^{\frac{(m-1)(n-1)}{4}} 2^{\frac{(m-1)(n-1)}{2}} \overset{\alpha}{\Pi}\overset{\beta}{\Pi} (\alpha - \beta).}
